\documentclass[11pt]{article}
\usepackage[utf8]{inputenc}
\usepackage[T1]{fontenc}
\usepackage{amsmath,amssymb,amsthm}
\usepackage{booktabs}
\usepackage{array}
\newcolumntype{L}[1]{>{\raggedright\arraybackslash}p{#1}}
\usepackage{graphicx}
% ---------------------------------------------------------------------------
% LISIBILITE — docs/articles/STYLE_ARTICLES.md §2.
% ---------------------------------------------------------------------------
\usepackage[a4paper,textwidth=13.8cm,textheight=22.4cm]{geometry}
\linespread{1.06}
\setlength{\parskip}{.55em plus .15em minus .1em}
\setlength{\parindent}{1.1em}
\emergencystretch=1.5em
\usepackage{hyperref}

% ============================================================================
% A4 — LE MARCHEUR. Chiffres : article/marcheur/MESURES.md, JAMAIS le journal.
% Chaque section porte en commentaire l'ANCRE (le fichier du depot qui la
% justifie). Regle d'or : une phrase = un objet du depot.
% ============================================================================

\newtheorem{definition}{Definition}
\newtheorem{proposition}{Proposition}
\theoremstyle{remark}\newtheorem{remark}{Remark}

\title{The Walk Across the Last Mile:\\
a Kernel-Guarded Search That Invents Its Own Waypoints}

\author{Karl Olivet\thanks{Independent researcher. Contact:
\texttt{karl.olivet@gmail.com}. Code and certified corpus:
\url{https://github.com/Karl-O/bourbaki-theorie-des-ensembles}; every figure
reported here was measured on 21 August 2026 and is recorded, with its
experiment, in \texttt{article/marcheur/MESURES.md}.}}

\date{Preprint --- August 2026}

\begin{document}
\maketitle

\begin{abstract}
%% ANCRE : MESURES.md §1 (la porte), §2 (mineur), §3 (oracle).
A companion paper argued that in mechanised proof search the last mile is
located, not crossed: a search that fails should name exactly what it lacks.
This paper reports the next step --- a \emph{walker} that crosses the stretch
the previous system could only locate. Given a goal that backward chaining
alone cannot close --- measured: on a four-element reassociation over a derived
operation $a \oplus b := (a+b)+1$, chaining exhausts its budget in 692\,s ---
the walker mines the goal's own subterms for repeated patterns, rediscovers
$\oplus$ without being told it exists, conjectures laws for it from an open
list of schemas, lets a numeric oracle kill the false ones in milliseconds,
has the kernel certify the true ones, and re-attempts the goal on the
\emph{compressed} pool of derived lemmas alone. The same goal then closes,
kernel-checked, in about 73\,s. Two measured facts frame the contribution.
First, compression is a \emph{replacement}, not an addition: re-attempting on
the derived lemma alone takes 72.6\,s, while adding it to the raw laws takes
962\,s --- a factor of 13, because new knowledge enlarges the search space
unless the old knowledge is set aside. Second, the walker's own first
component failed in a way worth reporting: ranking candidate patterns by size
selected only fragments of the huge $\tau$-developments (a triple-$\oplus$
term has $\approx$13{,}700 nodes) and missed $\oplus$ entirely; the fix ---
pairing subterms by root signature --- is itself a measured design law. As
everywhere in this series, the kernel guards every step: a bad proposal costs
a dead route, never a false theorem.
\end{abstract}

%% ============================================================================
\section{Introduction}
\label{sec:intro}
%% ANCRE : PLAN_ARTICLES.md (la porte, 10 aout) ; MESURES.md §1.
%% ============================================================================

The companion paper on failure reporting ended on a gate it refused to cross:
its editorial plan states that no walker should be published before it has
closed at least one goal that chaining alone cannot close. This paper begins
by crossing that gate, measured on both sides, and everything else in it is a
consequence of how the crossing was done.

\paragraph{The claim, in one paragraph.} A goal that defeats a bounded
backward search does not usually defeat it by depth alone. What is missing is
a \emph{waypoint}: an intermediate statement that, once certified and placed
in the pool, makes the remaining distance short. The walker's job is to invent
candidate waypoints from the goal's own structure, discard the false ones
cheaply, pay the kernel's price for the true ones, and then attempt the goal
from the waypoint rather than from the raw laws. None of this requires the
machine to be right: every candidate is judged --- first by a numeric oracle
that can only refute, then by the kernel that alone can certify. A wrong
proposal costs seconds; an unsound one is impossible.

%% A REMPLIR (EXP4) : le paragraphe de contributions avec le chiffre de la
%% marche de bout en bout, une fois exp4.out lu et colle dans MESURES.md §4.

%% ============================================================================
\section{Background: the bench, the budget, the gap}
\label{sec:background}
%% ANCRE : test_autonomie.py (v16-v18) ; MESURES.md de A2 ; besoin.py.
%% ============================================================================

We reuse the toy-algebra bench on which the failure-reporting paper measured
its rewriting organs. The base theory provides cardinal addition $a+b$ and its
two raw laws --- iterated associativity and commutativity --- as the only
facts in the pool. On top of it, a derived operation is defined and never
named anywhere in the repository:
\[
a \oplus b \;:=\; (a+b)+1 .
\]
The engine's budget is not a guess: the minimal raw rewriting chain for the
three-element associativity of $\oplus$ was measured at five steps, and five
is therefore the engine's cap (\texttt{max\_pas=5}); with a cap of three, that
same goal failed. All experiments below inherit this measured budget.

The goal that crosses the gate is the four-element reassociation
\[
B_4 \;:\;\; ((a \oplus b) \oplus c) \oplus d \;=\; a \oplus (b \oplus (c
\oplus d)),
\]
true, and out of reach of a five-step raw chain: each application of
$\oplus$-associativity costs about five raw steps, and $B_4$ needs two.

%% ============================================================================
\section{The gate, measured on both sides}
\label{sec:porte}
%% ANCRE : MESURES.md §1 (EXP1-EXP3).
%% ============================================================================

\begin{center}\small
\begin{tabular}{@{}L{0.42\linewidth}L{0.24\linewidth}r@{}}
\toprule
experiment & verdict & time \\
\midrule
chaining alone on $B_4$ (raw pool) & \textbf{fails}, one gap named & 692.5\,s \\
certifying the lemma $\oplus$-assoc & closed, kernel-checked & 4--7\,s \\
$B_4$ on the \emph{compressed} pool (lemma alone) & \textbf{closed}, 0 hypotheses & 72.6\,s \\
$B_4$ on the cumulative pool (raw $+$ lemma) & closed & 962.4\,s \\
false variant $F_4$ on the compressed pool & stays open (correct) & 79.1\,s \\
\bottomrule
\end{tabular}
\end{center}

Three readings, in decreasing order of importance.

\paragraph{The gate holds.} The same goal, from the same pool, under the same
budgets: chaining exhausts itself in 692\,s and returns a named gap; the
two-step walk --- certify the waypoint, then re-attempt from it --- closes it.
``Cannot close'' here is not an asymptotic claim; it is the measured behaviour
of the full 21-organ engine at its measured budget.

\paragraph{Compression is a replacement, not an addition.} The most useful
number in the table is the least intuitive one: adding the certified lemma to
the raw laws makes the search \emph{thirteen times slower} than replacing them
with it (962\,s against 72.6\,s). Every fact in the pool is a rewriting
candidate at every node; new knowledge enlarges the space unless the old
knowledge is set aside. The walker therefore re-attempts goals on the derived
lemmas \emph{alone} --- and says so in its journal, rather than silently.

\paragraph{Certification cost is noisy.} The same lemma certified twice took
3.8\,s and then 7.4\,s depending on cache state. We quote ``a few seconds''
and refuse a single figure.

%% ============================================================================
\section{The walker}
\label{sec:marcheur}
%% ANCRE : autonomie/marcheur.py (docstring de module).
%% ============================================================================

The walker's state is a pair (goal, pool of certified facts). One step:

\begin{enumerate}
\item \textbf{Sense.} Run the need organ. Closed: done. Open: a list of named
  gaps --- the failure object of the companion paper.
\item \textbf{Mine.} Extract repeated term patterns from the goal itself, by
  anti-unification with shared slots, ranked by the MDL gain
  $(\mathrm{occurrences}-1) \times \mathrm{size}$ --- the compression
  criterion of the notion-mining organ, transposed from formulas to terms.
\item \textbf{Conjecture.} Instantiate law schemas on each binary pattern:
  commutativity, associativity, idempotence. The schema list is \emph{open};
  nothing proves these three suffice, and the paper does not claim it.
\item \textbf{Refute cheaply.} A numeric oracle looks for small
  counterexamples. A false conjecture dies in milliseconds, never in minutes
  of kernel time. The oracle can only refute: no counterexample authorises
  spending, proves nothing.
\item \textbf{Certify, compress, retry.} Surviving conjectures go to the need
  organ over the raw pool; the kernel judges. Certified lemmas form the
  compressed pool on which the goal is re-attempted. No progress in a round:
  stop, and return the terminal gaps --- the walker, like the organ it wraps,
  fails by naming.
\end{enumerate}

The safety principle is unchanged since the first proposer organ: the walker
\emph{suggests}, the kernel \emph{judges}. No theorem object is built in the
walker's code; everything that enters the pool came out of the need organ,
hence out of the kernel.

%% ============================================================================
\section{The miner, including its first failure}
\label{sec:mineur}
%% ANCRE : MESURES.md §2 ; marcheur.py (miner_motifs, docstring).
%% ============================================================================

The miner answers the question the reader should be asking: \emph{who told
the machine about $\oplus$?} Nobody. The pattern is recovered from the goal.

\paragraph{Version 1 failed, and the failure is a design law.} The first
miner collected the goal's subterms and anti-unified the 24 largest. It found
nothing usable: in this kernel, terms are opaque $\tau$-developments, and a
triple-$\oplus$ term has ${\approx}13{,}700$ nodes, so the largest subterms
are all \emph{fragments of the developments' insides} --- and the small
instances $a \oplus b$, the ones that carry the pattern, never entered the
top-24. The top pattern was a three-slot context of MDL gain 13{,}737,
useless for conjecturing, and $\oplus$ was absent from the list. Ranking by
size selects noise.

\paragraph{Version 2 pairs by root signature.} Subterms are grouped by root
signature (tag, name, binder, arity) and each is anti-unified with its
size-neighbours \emph{within its group}: $O(n)$ pairs instead of $O(n^2)$,
and two instances of the same operation necessarily meet. Occurrences are
then counted over \emph{all} subterms, and candidates are validated by
reconstruction: a pattern applied to the pieces of its own generator must
rebuild it identically --- assemblages are hashable, so this is one $O(1)$
equality, never a structural guess. On $B_4$ the head pattern comes out with
6 occurrences and gain 4{,}045, and applying it to $(a,b)$ yields exactly
$\oplus(a,b)$: mining took 1.1\,s.

%% ============================================================================
\section{The oracle had a hole, and one layer was enough}
\label{sec:oracle}
%% ANCRE : MESURES.md §3 ; oracle_num.py (table, docstring).
%% ============================================================================

The refutation step exposed a measured blind spot. The oracle evaluates terms
by a reconstruction table --- numerals, their sums, their products --- because
terms here are opaque: there is nothing to descend into. But
$\mathrm{succ}(\text{sum})$ was in no entry, so \emph{every} conjecture about
a derived operation evaluated to ``unknown'': invisible, neither refutable
nor authorised.

One layer fixed it: for every term already in the table with value $v$, add
its successor with value $v+1$. After the extension, the idempotence schema
for $\oplus$ --- false --- dies in under a millisecond, before any kernel
call; commutativity and associativity survive the sweep (about a second each)
and proceed to certification. The limit is stated where the code is: one
layer only. A nested derived term such as $\mathrm{succ}(\mathrm{succ}(a+b)+c)$
is still outside the table, the full closure would cost $|\text{table}|^2$,
and on such conjectures the oracle stays silent --- it never affirms.

%% ============================================================================
\section{The walk, end to end}
\label{sec:marche}
%% ANCRE : MESURES.md §4 — A REMPLIR depuis exp4.out. NE RIEN ECRIRE AVANT.
%% ============================================================================

%% A REMPLIR (EXP4) : journal de marche reel (motifs, refutes, certifies,
%% FERME), duree totale, comparaison avec 692 s. Garde-fou F4 via la marche.

%% ============================================================================
\section{What the walk does not find}
\label{sec:limites-goldbach}
%% ANCRE : P6 — A MESURER (marcheur pointe sur un but Goldbach).
%% ============================================================================

%% A REMPLIR (P6) : le marcheur sur un but Goldbach ne ferme pas et rend ses
%% manques ; renvoi A3. NE RIEN ECRIRE AVANT LA MESURE.

%% ============================================================================
\section{Related work}
\label{sec:related}
%% A ECRIRE en dernier — regle 0-entree-(v?) : chaque citation verifiee a la
%% source dans ../references.bib.
%% ============================================================================

%% ============================================================================
\section{Limitations}
\label{sec:limitations}
%% ============================================================================

One bench. The walker has closed one family of goals over one derived
operation on one toy algebra. Nothing here claims generality: the paper's
claim is that the gate --- close something chaining cannot --- has been
crossed with every step measured, and that two of those measurements
(replacement-not-addition; size-ranking selects noise) are design laws that
survive the bench. The schema list is open. The walker does not learn: it
retains nothing between runs, and the title of the editorial plan for this
paper --- ``learning to propose'' --- has deliberately \emph{not} been kept,
because the code does not yet earn it.

%% ============================================================================
\section{Conclusion}
\label{sec:conclusion}
%% ============================================================================

%% A ECRIRE apres EXP4 et P6.

\bibliographystyle{alpha}
\bibliography{../references}

\end{document}
